\begin{frame}<1-3>[t]{新知探索}
\begin{campuscanvas}
% Source slide 16, objects 9; visible 2-3
\only<2-3|handout:1>{\node[anchor=north west,inner sep=0,text=black] at (70.000,150.000) {\begin{minipage}[t]{142.500mm}\raggedright\fontsize{16}{24.00}\selectfont \textbf{描述法表示集合的格式}\end{minipage}};}
% Source slide 16, objects 2; visible 3
\only<3|handout:1>{\node[anchor=north west,inner sep=0,text=black] at (70.000,235.000) {\begin{minipage}[t]{142.500mm}\raggedright\fontsize{11.5}{17.25}\selectfont 在一个大括号里写出集合中元素的共有属性。例如，集合 $\{2,4,6,8\}$ 可表示为 $\{\text{小于 10 的正偶数}\}$。\par 将“红马”看作集合，可表示为 $\{\text{红马}\}$，“赤兔马是红马”可表示为“赤兔马 $\in\{\text{红马}\}$”，这里的“是”相当于“$\in$”。\par 若用一句话描述不方便，通常先写出集合元素的一般属性或形式，再画一条竖线，竖线后列出元素要满足的条件。\par 例如，任何偶数都可表示为 $x=2k,\ k\in\mathbf Z$，所有偶数的集合为\[E=\{x\in\mathbf Z\mid x=2k,\ k\in\mathbf Z\}.\]\end{minipage}};}
\end{campuscanvas}
\end{frame}
